\uuid{t2sS}
\chapitre{Statistique}
\niveau{L2}
\module{Analyse de données}
\sousChapitre{Autre}
\titre{Variable quantitative à classer}
\theme{statistiques, tableur}
\auteur{}
\datecreate{2022-09-30}
\organisation{AMSCC}
\difficulte{}
\contenu{
\texte{Utiliser les données \href{https://github.com/smaxx73/Exercices/blob/main/data/ronfle.txt}{\texttt{ronfle.txt}} du dépôt (100 personnes, sept colonnes). Les colonnes numériques sont âge, poids, taille et alcool ; sexe, ronfle et taba sont des catégories. Les données sont séparées par des tabulations.}
\begin{enumerate}
\item \question{Importer le fichier dans un tableur et vérifier l'effectif ainsi que les types des colonnes.}
\reponse{Importer avec le séparateur tabulation, en conservant la première ligne comme en-tête. On obtient 100 observations, sans valeur manquante.}
\item \question{Classer la taille par intervalles de longueur 10 cm et tracer l'histogramme.}
\reponse{Avec les classes $[150;160[$, $[160;170[$, $[170;180[$, $[180;190[$, $[190;200[$ et $[200;210[$, les effectifs sont respectivement 4, 27, 10, 21, 37, 1. Leur somme vaut 100. Les classes étant de même largeur, on peut comparer leurs hauteurs directement.}
\item \question{Les âges sont-ils uniformément représentés ? Justifier graphiquement avec des classes de même largeur.}
\reponse{Pour $[20;30[$, $[30;40[$, $[40;50[$, $[50;60[$, $[60;70[$ et $[70;80[$, les effectifs sont 1, 14, 30, 24, 27, 4. L'histogramme est concentré entre 40 et 70 ans : l'échantillon ne présente pas une répartition uniforme. C'est un constat descriptif ; aucun test d'uniformité n'est demandé ici.}
\item \question{Tracer la taille en fonction du poids et commenter.}
\reponse{Utiliser un nuage de points : poids en abscisse, taille en ordonnée. L'association est croissante et forte ; la corrélation empirique vaut $r\simeq0{,}92697$. Cette corrélation descriptive n'établit pas de causalité.}
\item \question{Proposer un graphique faisant apparaître la répartition H/F selon des catégories de consommation d'alcool.}
\reponse{Un choix est de regrouper alcool en 0, de 1 à 4, puis 5 et plus. Les effectifs (F,H) sont respectivement $(23,19)$, $(2,31)$ et $(0,25)$ : 25 femmes et 75 hommes au total.
Utiliser des barres empilées pour les effectifs, ou des barres empilées à $100\%$ pour comparer les proportions H/F dans chaque catégorie. Indiquer explicitement les regroupements et le dénominateur des pourcentages. Le fichier ne précise pas l'unité de la variable alcool : ne pas en inventer une.}
\end{enumerate}
\indication{Pour les classes $[a;b[$, utiliser un comptage avec les conditions « supérieur ou égal à $a$ » et « strictement inférieur à $b$ ». Pour H/F selon alcool, utiliser un tableau croisé dynamique.}
}
